\section{Заключение}

В ходе работы были получены следующие результаты:

\begin{enumerate}
	\item Выведены уравнения движения перевернутого маятника с подвижной точкой опоры.
	\item Показано, что неуправляемый маятник неустойчив в верхнем положении.
	\item Для стратегии управления \(\ddot{u} = a\varphi + b\dot{\varphi}\) получены условия устойчивости: \(a > g,\; b > 0\).
	\item Для полной стабилизации (угол + точка опоры) предложено управление \(\ddot{u} = -a\varphi - b\dot{\varphi} - cu - d\dot{u}\).
	\item Сформулированы условия устойчивости на основе критерия Рауса-Гурвица.
	\item Проведены численные эксперименты, подтверждающие теоретические выводы.
\end{enumerate}

Таким образом, задача стабилизации перевернутого маятника успешно решена с использованием линейной обратной связи по состоянию.